\documentclass[12pt]{article}
\usepackage[a4paper,margin=1.5cm]{geometry}
\usepackage{amsmath}
\usepackage{amssymb}
\usepackage{bm}
\usepackage{graphicx}
\usepackage{float}

\setlength{\parindent}{0pt}
\setlength{\parskip}{0.4em}

\begin{document}

\section*{Table of Contents}
\textbf{1) Deriving the equations of motion} \par
\hspace{1cm} 1.1) Definitions \par
\hspace{1cm} 1.2) Kinematics \par
\hspace{1cm} 1.3) Dynamics \par
\hspace{1cm} 1.4) Euler-Lagrange equations \par
\hspace{1cm} 1.5) Converting the system of second order ODEs to a system of first order ODEs \par
\vspace{0.5cm}

\textbf{2) Exploring what makes the system ``chaotic''} \par
\hspace{1cm} 2.1) Time evolution of two trajectories \par
\hspace{1cm} 2.2) Characterizing the evolution of the perturbation vector \par
\vspace{0.5cm}

\textbf{3) Exploring global dynamics of the system} \par
\hspace{1cm} 3.1) State space reduction \par
\hspace{1cm} 3.2) Defining the goal \par
\hspace{1cm} 3.3) Finding a range for the initial values of \ensuremath{\theta,\omega_1} such that \ensuremath{\phi=0} and \ensuremath{\omega_2>0} is satisfied \par
\hspace{1.9cm} at \ensuremath{t=0} \par
\hspace{1cm} 3.4) Creating and visually analyzing the Poincare sections \par
\vspace{1cm}

\hrule
\vspace{1cm}

\section*{1) Deriving the equations of motion}

\subsection*{1.1) Definitions}

$l_1=\text{length of main pendulum rod}$

$l_2=\text{length of second pendulum rod}$

$\theta=\text{angle made by main pendulum rod with the vertical}$

$\phi=\text{angle made by the second pendulum rod with the vertical}$

$m_1=\text{mass of main pendulum}$

$m_2=\text{mass of second pendulum}$

\subsection*{1.2) Kinematics}

Position vector of the main pendulum:
\[
\mathbf{s}_1(\theta)=l_1(\sin\theta,\cos\theta)
\]

Velocity of the main pendulum:
\[
\mathbf{v}_1=l_1\dot{\theta}(\cos\theta,-\sin\theta)
\]
\[
|\mathbf{v}_1|^2=l_1^2\dot{\theta}^2
\]

Position vector of the second pendulum:
\[
\mathbf{s}_2(\theta,\phi)
=(l_1\sin\theta+l_2\sin\phi,l_1\cos\theta+l_2\cos\phi)
\]

Velocity vector of the second pendulum:
\[
\mathbf{v}_2
=(l_1\dot{\theta}\cos\theta+l_2\dot{\phi}\cos\phi,
-l_1\dot{\theta}\sin\theta-l_2\dot{\phi}\sin\phi)
\]
\[
|\mathbf{v}_2|^2
=l_1^2\dot{\theta}^2+l_2^2\dot{\phi}^2
+2l_1l_2\dot{\theta}\dot{\phi}\cos(\theta-\phi)
\]

\subsection*{1.3) Dynamics}

Kinetic energy of the system is:
\[
T=\frac{m_1}{2}l_1^2\dot{\theta}^2
+\frac{m_2}{2}\left[
l_1^2\dot{\theta}^2+l_2^2\dot{\phi}^2
+2l_1l_2\dot{\theta}\dot{\phi}\cos(\theta-\phi)
\right]
\]

Potential energy of the system is:
\[
U=-m_1gl_1\cos\theta-m_2g(l_1\cos\theta+l_2\cos\phi)
\]

Lagrangian of the system is:
\[
L=T-U
\]
\[
L=\frac{m_1}{2}l_1^2\dot{\theta}^2
+\frac{m_2}{2}\left[
l_1^2\dot{\theta}^2+l_2^2\dot{\phi}^2
+2l_1l_2\dot{\theta}\dot{\phi}\cos(\theta-\phi)
\right]
+m_1gl_1\cos\theta+m_2g(l_1\cos\theta+l_2\cos\phi)
\]

\subsection*{1.4) Euler-Lagrange equations}

Equation for $\theta$:
\[
\frac{\partial L}{\partial\theta}
=-m_2l_1l_2\dot{\theta}\dot{\phi}\sin(\theta-\phi)
-gl_1(m_1+m_2)\sin\theta
\]

\[
\frac{\partial L}{\partial\dot{\theta}}
=(m_1+m_2)l_1^2\dot{\theta}
+m_2l_1l_2\dot{\phi}\cos(\theta-\phi)
\]

\[
\frac{d}{dt}\left(\frac{\partial L}{\partial\dot{\theta}}\right)
=(m_1+m_2)l_1^2\ddot{\theta}
+m_2l_1l_2\ddot{\phi}\cos(\theta-\phi)
-m_2l_1l_2\dot{\theta}\dot{\phi}\sin(\theta-\phi)
+m_2l_1l_2\dot{\phi}^2\sin(\theta-\phi)
\]

\[
\frac{d}{dt}\left(\frac{\partial L}{\partial\dot{\theta}}\right)
=\frac{\partial L}{\partial\theta}
\]

\begin{equation}
\ddot{\theta}
+\frac{m_2}{m_1+m_2}\frac{l_2}{l_1}
\left(\ddot{\phi}\cos(\theta-\phi)
+\dot{\phi}^2\sin(\theta-\phi)\right)
+\frac{g}{l_1}\sin\theta=0
\tag{1}
\end{equation}

Equation for $\phi$:
\[
\frac{\partial L}{\partial\phi}
=m_2l_1l_2\dot{\theta}\dot{\phi}\sin(\theta-\phi)
-m_2gl_2\sin\phi
\]

\[
\frac{\partial L}{\partial\dot{\phi}}
=m_2l_2^2\dot{\phi}
+m_2l_1l_2\dot{\theta}\cos(\theta-\phi)
\]

\[
\frac{d}{dt}\left(\frac{\partial L}{\partial\dot{\phi}}\right)
=m_2l_2^2\ddot{\phi}
+m_2l_1l_2\ddot{\theta}\cos(\theta-\phi)
-m_2l_1l_2\dot{\theta}^2\sin(\theta-\phi)
+m_2l_1l_2\dot{\theta}\dot{\phi}\sin(\theta-\phi)
\]

\[
\frac{d}{dt}\left(\frac{\partial L}{\partial\dot{\phi}}\right)
=\frac{\partial L}{\partial\phi}
\]

\begin{equation}
\ddot{\phi}
+\frac{l_1}{l_2}
\left(\ddot{\theta}\cos(\theta-\phi)
-\dot{\theta}^2\sin(\theta-\phi)\right)
+\frac{g}{l_2}\sin\phi=0
\tag{2}
\end{equation}

\subsection*{1.5) Converting the system of second order ODEs to a system of first order ODEs}

Equations (1) and (2) can be represented in matrix form as:
\[
\begin{pmatrix}
1 & \mu\frac{l_2}{l_1}\cos(\Delta)\\
\frac{l_1}{l_2}\cos(\Delta) & 1
\end{pmatrix}
\begin{pmatrix}
\dot{\omega}_1\\
\dot{\omega}_2
\end{pmatrix}
=
\begin{pmatrix}
-\mu\frac{l_2}{l_1}\omega_2^2\sin(\Delta)
-\frac{g}{l_1}\sin\theta\\
\frac{l_1}{l_2}\omega_1^2\sin(\Delta)
-\frac{g}{l_2}\sin\phi
\end{pmatrix}
\]

Where $\frac{m_2}{m_1+m_2}=\mu$, $\theta-\phi=\Delta$, $\dot{\theta}=\omega_1$, $\dot{\phi}=\omega_2$

Which upon solving gives:
\[
\begin{pmatrix}
\dot{\omega}_1\\
\dot{\omega}_2
\end{pmatrix}
=
\begin{pmatrix}
\dfrac{
-\mu\frac{l_2}{l_1}\omega_2^2\sin(\Delta)
-\frac{g}{l_1}\sin\theta
-\mu\omega_1^2\cos(\Delta)\sin(\Delta)
+\mu\frac{g}{l_1}\cos(\Delta)\sin\phi
}{
1-\mu\cos^2(\Delta)
}
\\[10pt]
\dfrac{
\mu\omega_2^2\cos(\Delta)\sin(\Delta)
+\frac{g}{l_2}\cos(\Delta)\sin\theta
+\frac{l_1}{l_2}\omega_1^2\sin(\Delta)
-\frac{g}{l_2}\sin\phi
}{
1-\mu\cos^2(\Delta)
}
\end{pmatrix}
\]

Let $\mathbf{r}=(\theta,\omega_1,\phi,\omega_2)$ be the state vector of the system, then,

\[
\frac{d\mathbf{r}}{dt}
=(\dot{\theta},\dot{\omega}_1,\dot{\phi},\dot{\omega}_2)
\]

\[
\frac{d\mathbf{r}}{dt}
=(\omega_1,\dot{\omega}_1(\mathbf{r}(t)),\omega_2,\dot{\omega}_2(\mathbf{r}(t)))
\]

\[
\frac{d\mathbf{r}}{dt}=\mathbf{F}(\mathbf{r}(t))
\]
where
\[
\mathbf{F}(\mathbf{r}(t))
=(\omega_1,\dot{\omega}_1(\mathbf{r}(t)),\omega_2,\dot{\omega}_2(\mathbf{r}(t)))
\]

So we have successfully converted our system of two second order ODEs into four first order ODEs that can be numerically integrated using RK4.

\section*{2) Exploring what makes the system ``chaotic''}

\subsection*{2.1) Time evolution of two trajectories}

Numerically integrating the system of ODEs for the initial state $\mathbf{r}(0)=\left(\frac{2\pi}{3},0,\frac{\pi}{4},0\right)$ with $(m_1,l_1,m_2,l_2)=(5,3,7,2)$ and an infinitesimal perturbation of the initial state $\mathbf{r}'(0)=\mathbf{r}(0)+\mathbf{h}(0)$, where $\mathbf{h}(0)={10^{-6}}(1,0,1,0)$ allows us to plot angle vs time plots for the two trajectories. The plots for both $\theta$ vs $t$ and $\phi$ vs $t$ shows significant deviation between the two trajectories. This indicates sensitivity to initial conditions.

\begin{figure}[H]
    \centering
    \includegraphics[width=0.8\textwidth]{"theta vs t.png"}
    \label{fig:theta_vs_t}
\end{figure}

\begin{figure}[H]
    \centering
    \includegraphics[width=0.8\textwidth]{"phi vs t.png"}
    \label{fig:phi_vs_t}
\end{figure}

\subsection*{2.2) Characterizing the evolution of the perturbation vector}

Let $\mathbf{r}'(t)=\mathbf{r}(t)+\mathbf{h}(t)$ be another trajectory satisfying the ODE $\frac{d\mathbf{r}'(t)}{dt}=\mathbf{F}(\mathbf{r}'(t))$.

So,
\[
\frac{d(\mathbf{r}+\mathbf{h})}{dt}
=\mathbf{F}(\mathbf{r}(t)+\mathbf{h}(t))
\]

\[
\frac{d\mathbf{r}}{dt}+\frac{d\mathbf{h}}{dt}
=\mathbf{F}(\mathbf{r}(t)+\mathbf{h}(t))
\]

By performing a first order taylor expansion of $\mathbf{F}(\mathbf{r}(t)+\mathbf{h}(t))$ about $\mathbf{r}(t)$ we get:
\[
\frac{d\mathbf{r}}{dt}+\frac{d\mathbf{h}}{dt}
=\mathbf{F}(\mathbf{r}(t))
+[D\mathbf{F}(\mathbf{r}(t))]\mathbf{h}(t)
+O(|\mathbf{h}|^2)
\]

\[
\frac{d\mathbf{h}}{dt}
=[D\mathbf{F}(\mathbf{r}(t))]\mathbf{h}(t)
+O(|\mathbf{h}|^2)
\]

So upto first order our perturbation vector follows the ODE:
\[
\frac{d\mathbf{h}}{dt}
=[D\mathbf{F}(\mathbf{r}(t))]\mathbf{h}(t)
\]
where $[D\mathbf{F}(\mathbf{r}(t))]$ is the jacobian matrix of $\mathbf{F}(\mathbf{r}(t))$.

Numerically integrating both the peturbation vector ODE and the DP ODEs simultaneously for the initial state $\mathbf{r}(0)=\left(\frac{2\pi}{3},0,\frac{\pi}{4},0\right)$ with $\mathbf{h}(0)={10^{-6}}(1,0,1,0)$ and $(m_1,l_1,m_2,l_2)=(5,3,7,2)$ allows us to plot the logarithmic plot of $\ln\left(\frac{|\mathbf{h}(t)|}{|\mathbf{h}(0)|}\right)$ vs $t$. The obtained plot is quite jagged but is approximately linear.

\begin{figure}[H]
    \centering
    \includegraphics[width=0.8\textwidth]{"test for exponential divergence.png"}
    \label{fig:ln(|h(t)|/|h(0)|)_vs_t}
\end{figure}

The logarithmic plot is testing for exponential divergence of nearby trajectories in state space, for if:
\[
|\mathbf{h}(t)|\approx|\mathbf{h}(0)|e^{\lambda t}
\]
Then
\[
\ln\left(\frac{|\mathbf{h}(t)|}{|\mathbf{h}(0)|}\right)
\approx\lambda t
\]
should approximately be a linear graph with slope $\lambda$ where $\lambda$ is the maximal lyapunov exponent of the system.

If $|\mathbf{h}(t)|$ really does grow exponentially over time then for long term simulations, its magnitude would become arbitrarily large, causing the simulation to crash due to a floating point overflow error. To solve this issue and extract the maximal lyapunov exponent of the system, the benettin algorithm is used, which normalizes the perturbation vector at each step so its magnitude never becomes arbitrarily large while preserving its direction in state space. If $d_k$ is the magnitude of the perturbation vector at the kth time step, then the maximal lyapunov exponent is given by:
\[
\lambda=\frac{1}{T}\sum_{k}\ln(d_k)
\]
where $T$ is the time for which the simulation is run.

For the initial state $\mathbf{r}(0)=\left(\frac{2\pi}{3},0,\frac{\pi}{4},0\right)$ with $\mathbf{h}(0)=\frac{10^{-6}}{\sqrt{2}}(1,0,1,0)$ and $(m_1,l_1,m_2,l_2)=(5,3,7,2)$, the value of the maximal lyapunov exponent is found to be 0.7469. Since $\lambda>0$, nearby trajectories in state space divergence exponentially.

\begin{figure}[H]
    \centering
    \includegraphics[width=0.8\textwidth]{"convergence of the maximal lyapunov exponent.png"}
    \label{fig:lambda_vs_t}
\end{figure}

\section*{3) Exploring global dynamics of the system}

\subsection*{3.1) State space reduction}

The state space of the double pendulum is 4D. The total energy of the system is conserved and is given by:
\[
E=\frac{m_1l_1^2}{2}\omega_1^2
+\frac{m_2}{2}\left[
l_1^2\omega_1^2+l_2^2\omega_2^2
+2l_1l_2\omega_1\omega_2\cos(\theta-\phi)
\right]
-m_1gl_1\cos\theta
-m_2g[l_1\cos\theta+l_2\cos\phi]
\]

We can reduce the 4D state space to a 3D state space by using the total energy to find $\omega_2$ as a function of $\theta,\omega_1$ and $\phi$

\subsection*{3.2) Defining the goal}

For creating Poincare sections, we need a 2D state space. Because the double pendulum has only one constant of motion, the total energy, we need to introduce another physically meaningful constraint to be able to create a 2D section of the 4D state space. One physically meaningful constraint would be to record $\theta$ and $\omega_1$ values only when $\phi=0$ and $\omega_2>0$ i.e. study the motion of the main pendulum when the second pendulum crosses its lowermost point and has a positive angular velocity. To explore global behavior of the system under this constraint, we can integrate multiple trajectories at once, all of which satisfy $\phi=0$ and $\omega_2>0$ at $t=0$, and then store all the points $(\theta,\omega_1)$ when $\phi=0$ and $\omega_2>0$ for $t>0$ for each trajectory.

\subsection*{3.3) Finding a range for the initial values of \ensuremath{\theta,\omega_1} such that \ensuremath{\phi=0} and \ensuremath{\omega_2>0} is satisfied at \ensuremath{t=0}}

Imposing the constraint $\phi=0$ on the total energy at time $t=0$ we get:
\[
E=\frac{m_1l_1^2}{2}(\omega_1)_0^2
+\frac{m_2}{2}\left[
l_1^2(\omega_1)_0^2+l_2^2(\omega_2)_0^2
+2l_1l_2(\omega_1)_0(\omega_2)_0\cos\theta_0
\right]
-m_1gl_1\cos\theta_0
-m_2g[l_1\cos\theta_0+l_2]
\]

\[
E=\frac{(m_1+m_2)l_1^2}{2}(\omega_1)_0^2
+\frac{m_2l_2^2}{2}(\omega_2)_0^2
+m_2l_1l_2(\omega_1)_0(\omega_2)_0\cos\theta_0
-(m_1+m_2)gl_1\cos\theta_0-m_2gl_2
\]

Let $A=\frac{m_2l_2^2}{2}$, $B=m_2l_1l_2(\omega_1)_0\cos\theta_0$, $C=\frac{(m_1+m_2)l_1^2}{2}(\omega_1)_0^2-(m_1+m_2)gl_1\cos\theta_0-m_2gl_2$

So,
\[
E=A(\omega_2)_0^2+B(\omega_2)_0+C
\]

\[
0=A(\omega_2)_0^2+B(\omega_2)_0+C-E
\]

Thus,
\[
(\omega_2)_0=\frac{-B\pm\sqrt{B^2-4A(C-E)}}{2A}
\]

For $(\omega_2)_0\in\mathbb{R}$,
\[
B^2-4A(C-E)\geq0
\]

So,
\[
m_2^2l_1^2l_2^2(\omega_1)_0^2\cos^2\theta_0
-4\left(\frac{m_2l_2^2}{2}\right)
\left(
\frac{(m_1+m_2)l_1^2}{2}(\omega_1)_0^2
-(m_1+m_2)gl_1\cos\theta_0
-m_2gl_2-E
\right)\geq0
\]

Let $-(m_1+m_2)gl_1\cos\theta_0-m_2gl_2=V$ then,

\[
m_2^2l_1^2l_2^2(\omega_1)_0^2\cos^2\theta_0
-4\left(\frac{m_2l_2^2}{2}\right)
\left(
\frac{(m_1+m_2)l_1^2}{2}(\omega_1)_0^2+V-E
\right)\geq0
\]

Simplifying gives:
\[
(\omega_1)_0^2
\leq\frac{2(E-V)}
{l_1^2(m_2\sin^2\theta_0+m_1)}
\]

Thus,
\[
|(\omega_1)_0|
\leq
\sqrt{\frac{2(E-V)}
{l_1^2(m_2\sin^2\theta_0+m_1)}}
\]
\[
\Longrightarrow
-\sqrt{\frac{2(E-V)}
{l_1^2(m_2\sin^2\theta_0+m_1)}}
\leq(\omega_1)_0
\leq
\sqrt{\frac{2(E-V)}
{l_1^2(m_2\sin^2\theta_0+m_1)}}
\tag{1}
\]

Also since $2A>0$ always, for $\omega_2>0$ we need $-B+\sqrt{B^2-4A(C-E)}>0$.

Case 1) If $B<0$ then $-B>0$ and the inequality is satisfied automatically.

Case 2) If $B\geq0$ then $-B\leq0$:

Thus,
\[
\sqrt{B^2-4A(C-E)}>B
\]

So,
\[
B^2-4A(C-E)>B^2
\]
(since $B\geq0$)

Simplifying gives, $C<E$

So,
\[
\frac{(m_1+m_2)l_1^2}{2}(\omega_1)_0^2+V<E
\]

\[
(\omega_1)_0^2
<\frac{2(E-V)}{(m_1+m_2)l_1^2}
\]

Thus,
\[
|(\omega_1)_0|
<\sqrt{\frac{2(E-V)}{l_1^2(m_1+m_2)}}
\]
\[
\Longrightarrow
-\sqrt{\frac{2(E-V)}{l_1^2(m_1+m_2)}}
<(\omega_1)_0
<\sqrt{\frac{2(E-V)}{l_1^2(m_1+m_2)}}
\tag{2}
\]

Combining the two inequalities (1) and (2) for the case $B\geq0$ we get:
\[
\max\left(
-\sqrt{\frac{2(E-V)}{l_1^2(m_1+m_2)}},
-\sqrt{\frac{2(E-V)}
{l_1^2(m_2\sin^2\theta_0+m_1)}}
\right)
<(\omega_1)_0
<
\min\left(
\sqrt{\frac{2(E-V)}
{l_1^2(m_2\sin^2\theta_0+m_1)}},
\sqrt{\frac{2(E-V)}{l_1^2(m_1+m_2)}}
\right)
\]

Which can be simplified to:
\[
-\sqrt{\frac{2(E-V)}{l_1^2(m_1+m_2)}}
<(\omega_1)_0
<
\sqrt{\frac{2(E-V)}{l_1^2(m_1+m_2)}}
\]
since $m_2\sin^2\theta_0+m_1\leq m_1+m_2$ $\forall\theta$.

Also from both (1) and (2), for $\omega_1$ to be a real number, we need $E-V\geq0$

So, $E\geq V$

\[
\Longrightarrow
\frac{-E-m_2gl_2}{(m_1+m_2)gl_1}
\leq\cos\theta_0
\]

Let
\[
K=\frac{-E-m_2gl_2}{(m_1+m_2)gl_1}
\]

Case 1) If $K<-1$, then $\theta_0\in[-\pi,\pi]$.

Case 2) If $-1\leq K\leq1$, since cos is symmetric,
\[
-\cos^{-1}(K)\leq\theta_0\leq\cos^{-1}(K)
\]

Case 3) If $K>1$, then no physically valid $\theta_0$ exists.

So after choosing an energy value $E$ if:

1) For $K<-1$:
\[
-\pi\leq\theta_0\leq\pi
\]

For $-1\leq K\leq1$:
\[
-\cos^{-1}(K)\leq\theta_0\leq\cos^{-1}(K)
\]

2) For $B\geq0$:
\[
-\sqrt{\frac{2(E-V)}{l_1^2(m_1+m_2)}}
<(\omega_1)_0
<
\sqrt{\frac{2(E-V)}{l_1^2(m_1+m_2)}}
\]

For $B<0$:
\[
-\sqrt{\frac{2(E-V)}
{l_1^2(m_2\sin^2\theta_0+m_1)}}
\leq(\omega_1)_0
\leq
\sqrt{\frac{2(E-V)}
{l_1^2(m_2\sin^2\theta_0+m_1)}}
\]

Then $\phi=0$ and
\[
(\omega_2)_0=
\frac{-B\pm\sqrt{B^2-4A(C-E)}}{2A}>0
\]
is satisfied.

\subsection*{3.4) Creating and visually analyzing the Poincare sections}

Numerically integrating $\frac{d\mathbf{r}}{dt}=\mathbf{F}(\mathbf{r}(t))$ for 1000 trajectories, with $(m_1,l_1,m_2,l_2)=(5,3,7,2)$ and initial conditions as obtained in (3.3) and calculating the values of $(\theta,\omega_1,\omega_2(\theta,\omega_1))$ when $\phi=0$ and $\omega_2>0$ again for $t>0$ for each trajectory allows us to plot the Poincare section

\[
\left(\theta,\frac{\partial L}{\partial\dot{\theta}}\right)
\]

for various energy levels $E$. Giving each trajectory a color, we can see from the plots for $E\in\{-450\,\mathrm{J},-245\,\mathrm{J},0\,\mathrm{J},245\,\mathrm{J},490\,\mathrm{J}\}$, the presence of closed curves in the centre, with the curves becoming more and more defined as the energy increases. These closed curves can be identified with quasi-periodic orbits in phase space.

\begin{figure}[H]
    \centering
    \includegraphics[width=0.8\textwidth]{"-450 joules.png"}
    \label{fig:poincare_sec1}
\end{figure}

\begin{figure}[H]
    \centering
    \includegraphics[width=0.8\textwidth]{"-245 joules.png"}
    \label{fig:poincare_sec2}
\end{figure}

\begin{figure}[H]
    \centering
    \includegraphics[width=0.8\textwidth]{"0 joules.png"}
    \label{fig:poincare_sec3}
\end{figure}

\begin{figure}[H]
    \centering
    \includegraphics[width=0.8\textwidth]{"245 joules.png"}
    \label{fig:poincare_sec4}
\end{figure}

\begin{figure}[H]
    \centering
    \includegraphics[width=0.8\textwidth]{"490 joules.png"}
    \label{fig:poincare_sec5}
\end{figure}

\end{document}